\section[边界条件与扩张 TFT]{边界条件与扩张拓扑场论（3 月 2 日）}

\srcpage{29}

「这里有个双关，留给 Arun 当练习。」（注：我还没解出这道练习。）

注：我错过了周四的课。我会尽快看录像并更新笔记；在此期间，可以先看 Jackson 的笔记。

今天我们要考虑分歧（ramification），它在物理中对应于边界理论（boundary theory），或者对应于在奇点（singularity）上考虑算子。（TODO：我漏掉了一些物理例子，特别是螺线管（solenoid）与 Cartier 对偶。）

于是在等价 $\Pic(C)\simeq \mathrm{Loc}_1(C)$ 之下，我们有两种推广方式。可以在左边给线丛加上水平结构（level structure），即沿正余维的区域附加平凡化；在右边可以给局部系统加上奇点。而类域论的完整故事——把分歧也考虑进来——表明这个对应可以升级到包含这些推广。

在场论中，流形 $M$ 上有一个场空间（space of fields）$\cF(M)$，并有限制映射 $\pi_{\partial M}\colon\cF(M)\to\cF(\partial M)$。边界数据改变我们给 $N:=\partial M$ 附加向量空间的方式：具体地说，选取 $N$ 上的层 $\cE_N\to N$，用它写出另一个向量空间。例如，可以给 $M$ 附上向量空间 $\Gamma_M(\pi_{\partial M}^*\cE_N)$。

\begin{example}
  在拓扑 Maxwell 理论（topological Maxwell theory）中（A 侧），我们给 $M$ 附上的向量空间是 $H^*(\cC(M))$，即 $M$ 上 $U_1$-联络（connection）的空间的上同调。回忆至少在连通分支层面有 $\pi_0(\cC(M))\cong H^2(M;\Z)$。若考虑子流形 $N\subset M$ 并考虑边界数据，我们希望把它换成某种看起来像相对上同调（relative cohomology）的东西：$H^2(M,N;\Z)$。
\end{example}

\srcpage{30}

首先回忆 $H^2(M;\Z)$ 是怎么来的。主 $U_1$-丛上联络的空间可缩，而拓扑类型由 $BU_1\simeq K(\Z,2)$ 分类，所以 $\cC(M)\simeq \mathrm{Map}(M,K(\Z,2))$，即 $\cC(M)\simeq\mathrm{Map}(M,BU_1)$。现在，给定丛与联络，可以限制到 $N$，定义映射
\[
  \mathrm{Map}(M,K(\Z,2)) \longrightarrow \mathrm{Map}(N,K(\Z,2)) = \mathrm{Map}(N,BU_1).
\]
可以考虑的一个层是 $\cE:=i_*\C$（其中 $i\colon \mathrm{pt}\hookrightarrow N$ 是包含映射）；于是 $H^*(\cC(M);\pi_N^*\cE)$ 是「$M$ 上带 $N$ 上平凡化数据的线丛」的空间的上同调。所以这是 Dirichlet 边界条件（Dirichlet boundary condition）的一个例子。

在连通分支上，这恰好是 $H^2(M,N;\Z)$，即「映射 $f\colon M\to K(\Z,2)$ 连同限制 $f|_N$ 的零伦（null-homotopy）」的空间的连通分支群。这个空间是如下图表的拉回：
\[
  \begin{tikzcd}
    \arrow[r] \arrow[d] & \mathrm{Map}(M,K(\Z,2)) \arrow[d] \\
    \mathrm{pt} \arrow[r] & \mathrm{Map}(N,K(\Z,2)).
  \end{tikzcd}
\]

对偶地，可以考虑 Neumann 边界条件：设 $M_0$ 是 $M$ 中 $N$ 的管状邻域（tubular neighborhood）的补集，于是 $M_0$ 是带边流形。我们再次取 $\cF(\partial M_0)\leftarrow\cF(M_0)$ 上的常值层，并拉回（还是常值层），考虑 $H^*(\cF(M_0),\C)$。这是「可以在 $N$ 上有任意奇点的场」的空间的上同调。但也可以取摩天层的拉回，这等价于考虑 $M_0$ 上「在边界上带平凡化」的联络空间的上同调。你还可以在这里放更多一般的条件；边界上的每个层都是一台机器，给你体流形（bulk manifold）上一个新的向量空间。

我们过一会儿再回到这件事的数论方面，但首先来看扩张拓扑场论（extended topological field theory）！

回忆拓扑场论给一个闭 $n$-流形附上某个数（配分函数，partition function），可以启发式地把它想成场空间的体积。而给一个闭 $(n-1)$-流形 $N$，我们附上一个复向量空间（常常是 Hilbert 空间），把它想成 $M$ 上场空间上的泛函（functional）空间。我们得到这一点的方式是打破相对论对称性，把 $N$ 想成其圆柱中某一瞬间的截面。

现在我们再往前走一步：给定闭 $(n-2)$-流形 $P$，我们将指派一个范畴，常常是像 $\cF(P)$ 上的层范畴那样的东西。一般地，这个范畴是「线性的」，例如 Hom 空间是复向量空间；因此这是场的另一种线性化（linearization），只不过更加范畴化（categorified）。

传统上，这来自 Atiyah–Segal（TODO：此处应引文献）对拓扑场论的定义，即一个对称幺半函子（symmetric monoidal functor）
\[
  Z\colon \mathcal{B}ord_{n-1,n}\longrightarrow \Vect.
\]
这里 $\mathcal{B}ord_{n-1,n}$ 是这样一个范畴：对象是闭定向 $(n-1)$-流形，态射 $N_0\to N_1$ 是从 $N_0$ 到 $N_1$ 的协边（bordism）的微分同胚类\footnote{这个括号里的说明很重要，它保证复合在严格意义下结合。}。「对称幺半」意味着流形的不交并送到 $\Vect$ 中的张量积。复合是粘合协边。

想法是：从 $N_0$ 到 $N_1$ 的协边给出其边界上所附向量空间之间的线性映射 $Z(N_0)\to Z(N_1)$。而闭 $n$-流形 $M$ 是从 $\emptyset$ 到 $\emptyset$ 的协边，因此定义一个线性映射 $\C\to\C$；我们把它等同于一个复数，即 $M$ 的配分函数。

协边的视角告诉我们，可以把 $M$ 切成小块（具体说是协边），计算 $Z$ 在这些协边上的取值，再复合起来得到配分函数。物理上说，这表达了配分函数在时空中的局部性（locality）。

但态空间（state space）也应当是局部的！扩张 TFT 告诉我们如何把这个故事降低一维、升高一阶范畴。若闭 $(n-1)$-流形 $N$ 分解为协边的复合 $N=N_1\cup_P N_2$，我们希望表达
\[
  Z(N) = \langle Z(N_1), Z(N_2)\rangle_{Z(P)},
\]
其中配对要在某种范畴化的意义下理解，就像我们通过粘合计算配分函数一样。这个观点由 Freed、Lawrence 与 Dolan 等人阐述过。

\srcpage{31}

为了使之精确，我们构造一个协边 2-范畴 $\mathcal{B}ord_{n-2,n}$，提供如下数据。
\begin{itemize}
  \item \textbf{对象}：闭定向 $(n-2)$-流形。
  \item \textbf{态射}：紧定向协边，同前。
  \item \textbf{2-态射}：在 2-范畴中，态射之间现在有「2-态射」。这里允许协边之间的协边 $N_0$ 与 $N_1$，把它们想成带角流形（manifold with corners）。
\end{itemize}
对于 2-态射，$\partial N_0=\partial N_1$，所以不应把 $[0,1]\times[0,1]$ 想成带角流形的原型，而应想成眼形商（eye-shaped quotient）
\[
  ([0,1]\times[0,1])\big/\bigl((0,a)\sim(0,b),\ (1,a)\sim(1,b)\bigr).
\]

协边 2-范畴在不交并下仍是幺半的，尽管对称幺半 2-范畴的定义要麻烦一些。

\begin{definition}[两层 TFT / 一次扩张 TFT]
  两层 TFT（two-tier TFT），或一次扩张 TFT（once-extended TFT），是 2-范畴之间的对称幺半函子
  \[
    Z\colon \mathcal{B}ord_{n-2,n}\longrightarrow \cC,
  \]
  其中 $\cC$ 是某个对称幺半 2-范畴。
\end{definition}

此时 $\cC$ 该取什么并不显然，所以我们把它留作一个选择。不过一般地，我们希望它推广 Atiyah–Segal 的 TFT 定义，因此规定 $\cC$ 是 $\Vect$ 的一个 delooping\yizhu{delooping 指把一个幺半范畴实现为某个 2-范畴中对象的自同态范畴；这里要求 $\Omega\cC:=\End(\mathbf{1}_{\cC})$ 还原为 $\Vect$。中译尚无定名，保留原词。}{B（nLab "delooping"）}。也就是说，$\Omega\cC$（定义为 $\End(\mathbf{1}_{\cC})$）作为对称幺半范畴等价于 $\Vect$。

现在，若 $(n-1)$-流形 $N$ 如前分解为 $N_1\cup_P N_2$，就可以把 $N$ 写成 $\mathcal{B}ord_{n-2,n}$ 中态射的复合，从而至少原则上利用局部性计算 $Z(N)$。

要在实践中计算，应当选择 $\cC$。有多个好用的选择，但我们取 $\cC=\mathcal{C}at_\C$，即 $\C$-线性范畴的对称幺半 2-范畴。具体地说：
\begin{itemize}
  \item \textbf{对象}：（小）充实于 $\Vect_\C$ 的范畴。
  \item \textbf{态射}：尊重这种充实的函子。
  \item \textbf{2-态射}：尊重这种充实的自然变换。
\end{itemize}
果然，$\Omega\mathcal{C}at_\C\simeq\Vect$。

虽然我们没有把（拓扑）Maxwell 理论正式定义成 Atiyah–Segal TFT，从这个视角来想仍然有用：给定闭定向 2-流形 $\Sigma$，Maxwell 理论应当指派一个 $\C$-线性范畴；而若把闭 3-流形 $N$ 分解为 $N\cong N_1\amalg_\Sigma N_2$，就可以用这个分解计算指派给 $N$ 的向量空间。

阿贝尔对偶在余维 2 看起来是什么样？有几个例子值得研究（在 1 维及以下这是空话，所以从 2 维开始）。

\begin{example}
  在 2 维，我们有两个阿贝尔对偶的例子：一个里 $H^0(S^1;\Z)\leftrightarrow H^1(S^1;\C^\times)$，另一个里 $H^1(S^1;\Z)\leftrightarrow H^0(S^1;\C^\times)$。它们被解释为不同的场论；我们把它们扩张到余维 2。先看第一个：
  \begin{itemize}
    \item 在 A 侧（$H^0(S^1;\Z)$），场是到 $\Z$ 的映射。这也许不是最有趣的，但我们可以扩张：$\cF(\mathrm{pt})=\Z$，并通过取层来线性化。这里能取的层没有多少选择，最终 $Z(\mathrm{pt})=\mathcal{S}hv(\Z)$，即 $\Z$-分次向量空间的范畴。这是 $\C$-线性的，很好。
    \item 现在在 B 侧，对应于 $H^1(S^1;\C^\times)$。场是 $\C^\times$-局部系统，在 $\mathrm{pt}$ 上相当无聊：$\C^\times$-局部系统的同构类集合只有一个元素，而它上面的层不仅无聊，还与 A 侧得到的东西不匹配。为补救这一点，我们必须引入……叠\footnote{如果你对叠一无所知，那也没关系；下次我们会讨论关于它们的一些重要事情。}。局部系统的空间会增强为一个叠 $\mathcal{L}oc_{\C^\times}(\mathrm{pt})$，它编码局部系统有对称性这一事实。在这个具体情形里，局部系统只有一个同构类，但它有很多同构，整整一个 $\C^\times$。因此这个叠等价于 $B\C^\times=\mathrm{pt}/\C^\times$。我们要取它上面的层，这里指「点上的层，连同（唯一）茎上的 $\C^\times$-作用」。也就是说，B 侧场论给一个点指派的是 $\C^\times$ 的表示范畴。这同样是 $\C$-线性的。
  \end{itemize}
  对偶意味着我们最好能得到 $\Vect_\Z$ 与 $\Rep_{\C^\times}$ 之间的范畴等价，这由同型分解（isotypic decomposition）给出：若 $V$ 是 $\C^\times$-表示，它分裂为
  \[
    V = \bigoplus_{n\in\Z} V_n,
  \]

\srcpage{32}

  其中 $V_n\subset V$ 是 $z\in\C^\times$ 以 $z^n$ 作用的子空间。
\end{example}

这可以用 Fourier 的方式来看：把这两个范畴分别想成 $\C[\Z]$-模（A 侧）与 $\C[\C^\times]$-模（B 侧）的范畴。Fourier 变换需要稍微想一下余模（comodule），所以眼下不深入细节。

现在看另一个例子，$H^1(S^1;\Z)\leftrightarrow H^0(S^1;\C^\times)$。同样，A 侧涉及拓扑，B 侧涉及复几何。
\begin{itemize}
  \item A 侧用的是 $H^1(S^1;\Z)$ 上的函数。可以把它想成 $\Z$-值局部系统的同构类，但更方便的是映射到 $S^1=B\Z$，即周期函数。现在尝试扩张到一个点：场是到点的映射，于是 $\cF(\mathrm{pt})=S^1=B\Z$。取层时这可能不止一种意思，但从拓扑观点只有一种概念：局部常值层（locally constant sheaf），即 $\mathcal{L}oc_\Z(S^1)$。它们是 $\Z$ 的表示，或 $B\Z$ 上的层，这里 $\Z$ 来自 $\pi_1(S^1)$。于是最终得到 $\C[\Z]$-模，$\C[\Z]$ 上的代数结构来自卷积（convolution）。

  这东西有个花哨名字：$T^*S^1$ 的 wrapped Fukaya 范畴（wrapped Fukaya category）。一般地，对给定的辛流形（symplectic manifold）$X$，这是个难以研究的对象；但当 $X$ 是余切丛时，它可以被描述为局部系统。
  \item 在 B 侧，我们有 $H^0(S^1;\C^\times)$。场是到带离散拓扑的 $\C^\times$ 的映射（这一侧做的是代数几何，不是拓扑）。于是点上的场是 $\C^\times$；取层时，这里用代数几何，指的是 $\C^\times$ 上的拟凝聚层（quasicoherent sheaf）。这就是 $\C[\C^\times]$-模的范畴，代数结构是逐点乘法。Pontrjagin 对偶把带逐点乘法的 $\C[\C^\times]$ 与带卷积的 $\C[\Z]$ 等同起来，因此我们看到 $\mathcal{L}oc(S^1)\simeq \QCoh(\C^\times)$。

  这是同调镜像对称（homological mirror symmetry）的一个非常特殊的例子；一般地，它把 A 侧的 Fukaya 范畴与 B 侧的拟凝聚层范畴联系起来（在辛拓扑中，这就是 A 侧更偏拓扑的原因；因此我们在意层的代数几何概念）。
\end{itemize}

你可以在 3 维跑一个大致类似的游戏，但我们直接跳到关键例子：4 维 Maxwell 理论。在 A 侧，我们关心线丛的拓扑；在 B 侧，我们关心平坦线丛的代数几何。我们已经理解它在余维 1 如何运作：那里 $H^2(M;\Z)$ 与 $H^1(M;\C^\times)$ 之间有 Pontrjagin–Poincaré 对偶，其中 $M$ 是闭定向 3-流形；现在来看它给闭定向拓扑曲面 $\Sigma$ 附上什么。

在 A 侧，我们只有拓扑，场是带联络的线丛。我们通过取这个空间的 $\pi_0$ 得到 $H^2$。于是给 $\Sigma$ 附上场空间
\[
  \Pic(\Sigma) := \mathrm{Map}(\Sigma, BU_1).
\]
在 $\Sigma$ 上选一个复结构；于是有更熟悉的 $\Pic(\Sigma)$，即 $\Sigma$ 上全纯或代数线丛的空间，是代数几何中的对象。这个空间依赖于复结构，但其拓扑类型不依赖，并且事实上同伦等价于 $\mathrm{Map}(\Sigma,BU_1)$。它同样分解为
\[
  0 \longrightarrow \mathrm{Jac}(\Sigma) \longrightarrow \Pic(\Sigma) \xrightarrow{\deg} \Z \longrightarrow 0,
\]
其中 $\mathrm{Jac}(\Sigma)$ 是次数 $0$ 的线丛的空间。设 $\Sigma$ 连通且亏格为 $g$；则 $\mathrm{Jac}$ 的同伦型是实维数 $2g$ 的环面 $\C^g/\Z^{2g}$。

\begin{remark}
  线丛有自同构，因此 $\Pic(\Sigma)$ 与 $\mathrm{Jac}(\Sigma)$ 其实是叠，但目前不必操心这件事。
\end{remark}

上述短正合序列让人想起复几何中的指数正合序列（exponential exact sequence）：
\[
  0 \longrightarrow H^1(\Sigma;\cO_\Sigma)/H^1(\Sigma;\Z) \longrightarrow H^1(\Sigma;\cO_\Sigma^\times) \longrightarrow H^2(\Sigma;\Z).
\]
这些等同 $H^1(\Sigma;\cO_\Sigma^\times)\cong\Pic(\Sigma)$（同样忽略叠性）与 $H^2(\Sigma;\Z)\cong\Z$，连同连接同态与次数映射，意味着
\[
  \mathrm{Jac}(\Sigma) \cong H^1(\Sigma;\cO_\Sigma)/H^1(\Sigma;\Z).
\]
这把 Jacobian 等同于 $K(H^1(\Sigma;\Z),1)$，即以 $H^1(\Sigma;\Z)$ 为基本群的环面。这个环面也可以描述为 $S^1\otimes_\Z H^1(\Sigma;\Z)$。

\srcpage{33}

我们给 $\Sigma$ 附上的是 $\C$-线性范畴 $\mathrm{Loc}(\Pic(\Sigma))$\footnote{叠性不影响这个阿贝尔范畴，但会影响导出范畴。}。周四我们会看到它与 B 侧理论给 $\Sigma$ 附上的范畴有什么关系。
